\chapter{Discussion and Conclusions}\label{ch:conclusions}

\section{Answers to the research questions}\label{sec:conclusions-findings}

The derivations are in the chapters named; the answers are stated here without
repeating them.

\paragraph{RQ1: the joint score and its couplings.} Linear and known in closed
form, $\score(x,t) = -\Qt x$ with $\Qt = (e^{-2t}\Sigz + \Deltat I)^{-1}$. For
$\alpha \neq 0$ its
couplings have a closed-form lifecycle: exactly tridiagonal at $t=0$, filling
one further diagonal per order in $t$ with every off-band entry provably
non-zero at every $t > 0$, maximally coupled at intermediate times in the
Frobenius measure used to track it,
returning to the identity at rate $e^{-2t}$. A single bulk quantity
$q(\alpha,t) \in (0,|\alpha|]$ sets the posterior correlation length. This is a
consequence of the model's linearity and not of the chain structure alone
(Chapter~\ref{ch:gaussian}).

\paragraph{RQ2: belief propagation and the Kalman smoother.} They are the same
computation. The channel corrupts each coordinate independently, so every
observation enters through a unary factor, the posterior factor graph stays a
tree, and one forward and one backward sweep give all smoothing marginals. On
the Gaussian chain the messages close algebraically and the recursion is the
Kalman filter followed by RTS smoothing, update by update; the two routes agree
to $10^{-15}$--$10^{-14}$ across the tested grid, reproduced by an independently
written implementation. Tree exactness holds for any innovation law; closure in
a finite-dimensional family holds only because the model is Gaussian
(Chapter~\ref{ch:gaussian-bp}).

\paragraph{RQ3: what Gaussian messages cost.} Median relative score error
$0.27$ at $t=0.05$ on Laplace chains at $\alpha = 0.85$, below $10^{-2}$ by
$t \approx 1$ and
below $10^{-3}$ by $t \approx 2.1$, with direction preserved far better than
magnitude. Theorem~\ref{thm:l-lmmse} says what that number is: moment-projected
sum--product returns the LMMSE estimator of the covariance-matched Gaussian
model, so the measurement is full posterior inference against the best
second-order linear estimator. The \emph{estimator} is therefore fixed by the
model class --- the same for every innovation law with matching first two
moments --- while the measured \emph{error} against the true posterior is
law-dependent, as the innovation sweep shows: at $\alpha = 0.5$ and
$t = 0.05$ it spans $0.07$ for a mildly bimodal
mixture to $0.94$ for a strongly bimodal one at the same second moments
(Chapter~\ref{ch:laplace}).

\paragraph{RQ4: how fast remote evidence stops mattering.} In the bulk of a
long chain the influence of
evidence at distance $d$ decays exactly as $q(\alpha,t)^d$ --- an
infinite-chain statement, approximate near the boundaries of a finite one ---
and the error of the
radius-$r$ windowed estimator is measured to decay at a matching geometric
rate over the tested radii. These stay two claims: windowing changes the
recursion's
boundary conditions in a way the influence argument does not cover, so the
second is a numerical law over the tested range. At equal locality budget,
truncating the inference beats truncating the score matrix at small $t$
($12.3\%$ against $21.0\%$ at $t = 0.05$) (Chapter~\ref{ch:gaussian-bp}).

\paragraph{RQ5: recovering an unknown transition.} Fisher's identity turns one
forward--backward pass, at each current parameter value, into the exact
gradient of the marginal log-likelihood
of the functional model, so the E-step is structurally exact for any
innovation law and nothing is
differentiated through the recursion; the implementation evaluates it on a
validated grid, and recovery is demonstrated numerically over the studied
regime rather than proved asymptotically. The checks are part of the answer, not a
caveat on it: the channel acts unevenly across coordinates, and the fitted
innovation shape settles at a median of $\convshapemed$ updates against
$\convrhomed$ for the correlation. A stopping rule read off the correlation
trace therefore certifies a kernel that has not converged in the coordinate
every non-Gaussian claim depends on, so a fixed iteration budget compares
convergence rates and not estimators
(Chapters~\ref{ch:em-method} and~\ref{ch:em-kernel}).

\paragraph{RQ6: what supplying the Markov assumption is worth.} At an equal
number of clean trajectories, $\ratiolo$--$\ratiohi\times$ lower
schedule-pooled relative
score error than networks
trained by denoising score matching, across \nsizesused{} sizes; against a
baseline given locality and weight sharing, $\structratiolo$--$\structratiohi\times$.
The factor is quoted under the aggregation the text defines --- squared error
and reference energy pooled over the schedule per seed, ratio paired within
seed --- and Appendix~\ref{app:aggregation} shows the alternatives, of which
the retired mean of per-level ratios is systematically the largest.
Three scope conditions travel with those numbers. The arms are asymmetric in
supplied information by construction, and in both directions: EM--BP is given
the Markov factorisation, while every network arm is given the clean
trajectories and redraws the corruption at each gradient step, where EM--BP
sees one frozen noised realisation and never the clean signal. So this
measures what a correct structural assumption is worth against the standard
training recipe for the alternative, per clean trajectory, and it does not
rank learning algorithms at a matched information interface. The gap between the two ratios is not a decomposition into
architecture and estimation terms, because architecture, capacity, sharing and
optimisation change together when the baseline is swapped. And both are ratios
of errors at fixed data, not of the data each method needs for fixed error; the
curves do not cross in the measured range, so no data-equivalence factor follows
(Chapter~\ref{ch:em-results}).

\section{What is new}

\paragraph{The exact structure of a chain diffusion score.} The joint score of a
noised Markov chain is characterised completely for the Gaussian case: its
closed form, its eigenbasis at every diffusion time, the law by which its
couplings fill and then decay, and the single parameter that controls posterior
influence range. The identification of chain BP with the Kalman--RTS recursion
makes the two classical literatures one object rather than two analogous ones.

\paragraph{An exact reading of what Gaussian closure computes.} Where previous
practice measured the error of moment-projected message passing as the error of
an approximation, Theorem~\ref{thm:l-lmmse} identifies the estimator it returns.
The measured Laplace gap is therefore the price of discarding everything beyond
second moments --- a property of the model class and not of one
implementation. The estimator is shared by every innovation law with the same
first two moments; the size of the gap is each law's own.
The validated grid reference that makes the measurement possible is itself a
contribution, since the functional recursion has no closed form.

\paragraph{Transition learning, and the measured value of structure.} The
transition law is estimated from noised sequences alone by exact chain inference
inside EM, and the resulting estimator is compared against trained networks at
equal data under a protocol that selects both arms on held-out validation. Two
boundaries of the advantage are mapped rather than asserted: mixture capacity
beyond a single Gaussian innovation buys nothing for this estimator under the
grid and cap implemented here, and a chain-structured estimator partly absorbs a rank-one violation of the
Markov assumption through its refitted local parameter while failing on a
long-range one at a coupling strength of roughly one tenth. The convergence
measurement that separates rate from estimator is a methodological result in its
own right, and it is what withdrew an earlier reading of the generative
comparison.

\section{Limitations}

\begin{enumerate}
  \item \textbf{Model class.} The exact solutions concern a one-dimensional
  state per frame and linear AR(1) dynamics, extended to a two-dimensional
  nonlinear state only in Chapter~\ref{ch:ring}. The locality laws are proved
  for the Gaussian bulk; their non-Gaussian counterparts are only structurally
  delimited, and no low-rank structure is asserted for the Laplace
  $\nsites = 2$ posterior Hessian beyond the narrow regime where it holds.

  \item \textbf{The numerical reference.} The non-Gaussian results are
  numerical. Grid BP is a reference only within its validated resolution regime
  ($\sqrt{2t} \gtrsim 3h$ on the grid spacing $h$), and the Laplace closure
  error is a measured curve rather than a formula.

  \item \textbf{Optimisation and finite seeds.} The learning results are
  statistical rather than exact: estimates over a finite number of seeds,
  reported with paired uncertainty, several resolved at only one of the two
  training sizes tested. The headline factors are properties of a stated
  aggregation --- the schedule-pooled paired ratio --- and move by tens of
  percent under defensible alternatives (Appendix~\ref{app:aggregation}), so
  no exact multiple should be carried out of context. Every capacity cell
  above $C = \capcappedfrom$ stops at
  an iteration cap rather than at tolerance, so a statistical disadvantage of
  larger mixtures is not separated from an optimisation or a discretisation one.
  The convergence chapter's fitted slopes are coordinate-wise contraction
  factors, not Jacobian eigenvalues, and its sample-size table is exploratory.
  The misspecification study probes exactly two violations of the Markov
  assumption, with a scalar that is an exact KL only in the Gaussian arms and
  a covariance proxy in the Laplace ones; it establishes that the two
  mechanisms are not interchangeable in effect, not
  a general theory of how much non-Markov structure a chain estimator tolerates.

  \item \textbf{Generative fidelity, and transfer.} Whether pointwise accuracy
  and generative fidelity can move in opposite directions within one estimator
  is unresolved here: the comparison that addresses it evaluates procedures
  whose fitted shape coordinates are unequally settled, on the single coordinate
  the generative metric scores, so its reading is withdrawn rather than
  reversed. More broadly, these are properties of one controlled model. Whether
  trained score networks on real sequence data inherit them is untested here.
  The image and video extensions explored in the underlying research package are
  outside the scope of this thesis and are not reported.
\end{enumerate}

\section{Future work}

Four directions, in the order in which they would most change what this thesis
can claim.

\emph{The reverse-generation comparison, with converged fits.} The withdrawn
study integrated the reverse SDE under the estimated and the trained scores, but
on fits the convergence analysis shows to be unequally settled. Selecting both
arms on a convergence criterion rather than a fixed iteration count would answer
the question the experiment was built for, and close the one question this
thesis leaves open.

\emph{A grid-by-capacity design.} The capacity result leaves two mechanisms
entangled. A finer grid at fixed capacity would separate fits falling below the
grid's resolution from genuine overfitting, since only the second survives mesh
refinement; a budget ladder at fixed capacity would do the same for the
iteration cap.

\emph{Richer departures from the Markov assumption.} Two violations at fixed
strengths are probed here. A paired design across strengths, and further
violation types, would test whether the finding that the two mechanisms are not
interchangeable generalises. Discrete chains are untouched: with a finite
alphabet the messages are finite vectors and inference is exact with no closure
at all, giving a second solvable benchmark alongside the Gaussian one.

\emph{Spatial as well as temporal structure.} Whether the structural results
survive longer sequences, other innovation laws, and data with spatial
structure is the question that decides how far any of this carries beyond the
solvable setting. The reading of U-Nets as belief propagation
(Section~\ref{sec:sm-unet}) suggests where to look: on a hierarchical rather
than a chain graph, the same inference argument should apply, with the
tree-exactness that makes it work here replaced by an approximation whose cost
would have to be measured the way Chapter~\ref{ch:laplace} measures the
Gaussian closure.
